\documentclass{article}
\usepackage{graphicx} % Required for inserting images
\usepackage{amsmath}
\usepackage{amssymb}


\title{Analisi Matematica a.a.2026/27
seconda attività didattica}
\date{Novembre 2026}

\begin{document}

\maketitle

\paragraph{Definizione di \(f\)}

\[
f(x,y) = 
\begin{cases}
    \dfrac{|x+y||x-y-4|}{\sqrt{x^2 + y^2}}, & \text{se } (x,y) \neq (0,0) \\[6pt]
    0, & \text{se } (x,y) = (0,0)
\end{cases}
\]

Si ha che:

\[
\operatorname{Dom}f := \{(x,y) \in \mathbb{R}^2 \}
\]

Verifichiamo la continuità di \(f\) in \((0,0)\). Per fare ciò, calcoliamo il limite di \(f(x,y)\) quando \((x,y) \to (0,0)\).

\[
\lim_{(x,y) \to (0,0)} f(x,y) = \lim_{(x,y) \to (0,0)} \frac{|x+y||x-y-4|}{\sqrt{x^2 + y^2}}
\]

Il limite è una forma indeterminata del tipo \(\frac{0}{0}\). Per risolvere questo limite,consideriamo il percorso lungo la retta \(y = mx\):

Sostituendo:

\[
\lim_{(x,mx) \to (0,0)} \frac{|x+mx||x-mx-4|}{\sqrt{x^2 + m^2x^2}} = \lim_{x \to 0} \frac{|x||1 + m||x(1-m) - 4|}{|x|\sqrt{1+m^2}} =
\]
\[
= \frac{|1+m|}{\sqrt{1+m^2}} \lim_{x \to 0} |x(1-m) - 4| =\frac{|1+m|}{\sqrt{1+m^2}} \cdot 4
\]

Poichè il limite dipende da \(m\) non esiste. Pertanto, \(f\) non è continua in \((0,0)\).

\paragraph{Studio della differenziabilità in \((2,-2)\)}

Dalla definizione di differenziabilità, \(f\) è differenziabile in \((2,-2)\) se esistono \(\alpha, \beta \in \mathbb{R}\), tali che

\[
\lim_{(x,y) \to (2,-2)} \frac{f(x,y) - f(2, -2) - \alpha(x-2) - \beta(y+2)}{\sqrt{{(x-2)}^2 + {(y+2)}^2}} = 0
\]

Se \(f\) è differenziabile in \((2,-2)\), allora esistono le derivate parziali di \(f\) in \((2,-2)\) e si ha che \(\alpha = \frac{\partial f}{\partial x}(2,-2)\) e \(\beta = \frac{\partial f}{\partial y}(2,-2)\).

\[
\frac{\partial f}{\partial x}(2,-2) = \lim_{h \to 0} \frac{f(2+h,-2) - f(2,-2)}{h} = \lim_{h \to 0} \frac{f(2+h,-2) - 0}{h} = 
\]
\[
= \lim_{h \to 0} \frac{\frac{h^2}{\sqrt{{(2+h)}^2 + {(-2)}^2}}}{h}  = 0
\]


\[
\frac{\partial f}{\partial y}(2, -2) = \lim_{k \to 0} \frac{f(2,-2+k) - f(2,-2)}{k} = \lim_{k \to 0} \frac{f(2,-2+k) - 0}{k} = 
\]
\[
= \lim_{k \to 0} \frac{\frac{k^2}{\sqrt{{(2)}^2 + {(-2+k)}^2}}}{k}  = 0
\]

Poiche \(f(2,-2) = 0\), si ha che \(f\) è differenziabile in \((2,-2)\) se e solo se 

\[
\lim_{(x,y) \to (2,-2)} \frac{f(x,y)}{\sqrt{{(x-2)}^2 + {(y+2)}^2}} = 0
\]

Da cui segue che \(f\) è differenziabile in \((2,-2)\) se e solo se

\[
\lim_{(x,y) \to (2,-2)} \frac{\frac{|x+y||x-y-4|}{\sqrt{x^2 + y^2}}}{\sqrt{{(x-2)}^2 + {(y+2)}^2}} = \lim_{(x,y) \to (2,-2)} \frac{|x+y||x-y-4|}{\sqrt{x^2 + y^2} \cdot \sqrt{{(x - 2)}^2 + {(y + 2)}^2}} = 0
\]

il limite è una forma indeterminata del tipo \(\frac{0}{0}\). 

Utilizziamo le coordinate polari per calcolare il limite. Poniamo \(x = 2 + \rho\cos(\theta)\) e \(y = -2 + \rho\sin(\theta)\).

\[
\frac{|x+y||x-y-4|}{\sqrt{x^2 + y^2} \cdot \sqrt{{(x - 2)}^2 + {(y + 2)}^2}} = \frac{|(2 + \rho\cos(\theta)) + (-2 + \rho\sin(\theta))|| (2 + \rho\cos(\theta)) - (-2 + \rho\sin(\theta)) - 4|}{\sqrt{(2 + \rho\cos(\theta))^2 + (-2 + \rho\sin(\theta))^2} \cdot \sqrt{{(\rho\cos(\theta))}^2 + {(\rho\sin(\theta))}^2}}
\] 

\[
= \frac{\rho^2|(\cos(\theta) + \sin(\theta))(\cos(\theta) - \sin(\theta))|}{\rho\sqrt{{(2+\rho\cos(\theta)}^2 + {(-2+\rho\sin(\theta))}^2}} = \frac{\rho|(\cos(\theta) + \sin(\theta))(\cos(\theta) - \sin(\theta))|}{\sqrt{{(2+\rho\cos(\theta))}^2 + {(-2+\rho\sin(\theta))}^2}}
\]

\[
 = \rho \frac{|\cos(2\theta)|}{\sqrt{{(2+\rho\cos(\theta))}^2 + {(-2+\rho\sin(\theta))}^2}}
\]

Non è possibile esplicitare \(g(\rho)\) e \(h(\theta)\).

Tuttavia, si ha chè 

\[
0 \leq \frac{\rho|\cos(2\theta)|}{\sqrt{{(2+\rho\cos(\theta))}^2 + {(-2+\rho\sin(\theta))}^2}} \leq \frac{\rho}{\sqrt{8} - \rho} \to 0 \; \qquad (\rho \to 0^+)
\]

Pertanto, per il teorema del confronto, si ha che

\[
\lim_{(x,y) \to (2,-2)} \frac{|x+y||x-y-4|}{\sqrt{x^2 + y^2} \cdot \sqrt{{(x - 2)}^2 + {(y + 2)}^2}} = 0
\]

\(f\) è quindi differenziabile in \((2,-2)\) e le derivate parziali sono entrambe nulle.

\paragraph{Studio della differenziabilità in \((1, 1)\)}

Notiamo che \(x+y = 2 \neq 0\), \(x-y-4 = -4 \neq 0\) e \(x^2 + y^2 = 2 \neq 0\) in \((1,1)\)

Il che implica che in un intorno sufficientemente piccolo di \(1,1\) i segni delle espressioni dentro i valori assoluti non cambiano.

Pertanto, in un intorno di \((1,1)\), \(f\) coincide con la funzione 

\[
g(x,y) = \frac{(x+y)(-x+y+4)}{\sqrt{x^2 + y^2}} = \frac{-x^2 + y^2 + 4x + 4y}{\sqrt{x^2 + y^2}}
\]

Calcoliamo le derivate parziali di \(g\):

\[
\frac{\partial g}{\partial x} = \frac{(-2x+4)\sqrt{x^2 + y^2} - (-x^2 + y^2 + 4x + 4y)\cdot \frac{x}{\sqrt{x^2 + y^2}}}{x^2+y^2}
\]

\[
\frac{\partial g}{\partial y} = \frac{(2y+4)\sqrt{x^2 + y^2} - (-x^2 + y^2 + 4x + 4y)\cdot \frac{y}{\sqrt{x^2 + y^2}}}{x^2+y^2}
\]

Poichè \(\frac{\partial g}{\partial x}\) e \(\frac{\partial g}{\partial y}\) sono continue in \((1,1)\), segue che \(g\) è differenziabile in \((1,1)\) per il teorema del differenziale totale e quindi \(f\) è differenziabile in \((1,1)\) per il teorema del differenziale totale.

\paragraph{Studio della differenziabilità in \((0,0)\)} Poiché \(f\) non è continua in \((0,0)\), segue che \(f\) non è differenziabile in \((0,0)\).

\paragraph{Verifica dell'esistenza e eventuale calcolo di \(\partial_v f(1,1), v = \begin{pmatrix} \frac{1}{2} \\ -\frac{\sqrt{3}}{2}\end{pmatrix}\)}

Verifichiamo che \(v\) sia un vettore unitario:

\[
\langle v, v \rangle = \left(\frac{1}{2}\right)^2 + \left(-\frac{\sqrt{3}}{2}\right)^2 = \frac{1}{4} + \frac{3}{4} = 1 \implies ||v|| = 1
\]

Poiché \(f\) è differenziabile in \((1,1)\), segue che \(f\) è derivabile in ogni direzione in \((1,1)\) e quindi esiste la derivata direzionale di \(f\) in \((1,1)\) lungo \(v\).

Inoltre come conseguenza del teorema della derivata composta, si ha che 

\[
\partial_v f(1,1) = \langle \nabla_f(1,1), v \rangle
\]

Determiniamo \(\nabla_f(1,1)\):

\[
\frac{\partial f}{\partial x}(1,1) =  \frac{(-2+4)\sqrt{2} - (-1 + 1 + 4 + 4)\cdot \frac{1}{\sqrt{2}}}{2} = \frac{2\sqrt{2} - 8\cdot \frac{1}{\sqrt{2}}}{2} = \frac{2\sqrt{2} - 4\sqrt{2}}{2} = -\sqrt{2}
\]

\[
\frac{\partial f}{\partial y}(1,1) =  \frac{(2+4)\sqrt{2} - (-1 + 1 + 4 + 4)\cdot \frac{1}{\sqrt{2}}}{2} = \frac{6\sqrt{2} - 8\cdot \frac{1}{\sqrt{2}}}{2} = \frac{6\sqrt{2} - 4\sqrt{2}}{2} = \sqrt{2}
\]


Segue che \(\nabla_f(1,1) = \begin{pmatrix} -\sqrt{2} \\ \sqrt{2} \end{pmatrix}\), il che implica che

\[
\partial_v f(1,1) = \langle \nabla_f(1,1), v \rangle = \left\langle \begin{pmatrix} -\sqrt{2} \\ \sqrt{2} \end{pmatrix}, \begin{pmatrix} \frac{1}{2} \\ -\frac{\sqrt{3}}{2}\end{pmatrix} \right\rangle = -\frac{\sqrt{2}}{2} - \frac{\sqrt{6}}{2} = -\frac{\sqrt{2} + \sqrt{6}}{2}
\]

\paragraph{Verifica dell'esistenza e eventuale calcolo di \(\partial_v f(0,0), v = \begin{pmatrix} -\frac{1}{\sqrt{2}} \\ \frac{1}{\sqrt{2}}\end{pmatrix}\)}

Verifichiamo che \(v\) sia un vettore unitario:
\[
\langle v, v \rangle = \left(-\frac{1}{\sqrt{2}}\right)^2 + \left(\frac{1}{\sqrt{2}}\right)^2 = \frac{1}{2} + \frac{1}{2} = 1 \implies ||v|| = 1
\]

Poiché \(f\) non è differenziabile in \((0,0)\), non possiamo concludere che \(f\) sia derivabile in ogni direzione in \((0,0)\). Tuttavia, possiamo verificare se la derivata direzionale di \(f\) in \((0,0)\) lungo \(v\) esiste.

\[
\partial_v f(0,0) = \lim_{h \to 0} \frac{f(0 + h \cdot (-\frac{1}{\sqrt{2}}), 0 + h \cdot \frac{1}{\sqrt{2}}) - f(0,0)}{h} = \lim_{h \to 0} \frac{\frac{|-\frac{1}{\sqrt{2}}h + \frac{1}{\sqrt{2}}h||-\frac{1}{\sqrt{2}}h -\frac{1}{\sqrt{2}}h - 4|}{\sqrt{(-\frac{1}{\sqrt{2}}h)^2 + (\frac{1}{\sqrt{2}}h)^2}} - 0}{h} = 0
\]

\paragraph{Verificare l'esistenza e eventualmente determinare l'equazione del piano tangente al grafico di \(f\) nel punto \(1,1, f(1,1)\)}

Poiché \(f\) è differenziabile in \((1,1)\), segue che il piano tangente al grafico di \(f\) nel punto \((1,1,f(1,1))\) esiste ed è dato da

\[
z = f(1,1) + \frac{\partial f}{\partial x}(1,1)(x-1) + \frac{\partial f}{\partial y}(1,1)(y-1) = 4\sqrt{2} - \sqrt{2}(x-1) + \sqrt{2}(y-1) 
\]

\[
= 4\sqrt{2} - \sqrt{2}x + \sqrt{2} + \sqrt{2}y - \sqrt{2} = 4\sqrt{2} - \sqrt{2}x + \sqrt{2}y
\]

\end{document}
